% Part: methods
% Chapter: proofs
% Section: proof-by-contradiction

\documentclass[../../../include/open-logic-section]{subfiles}

\begin{document}

\olfileid[hi]{mth}{prf}{con}

\olsection{विरोधाभास द्वारा प्रमाण}

पहली स्थिति में विरोधाभास द्वारा प्रमाण ऐसा अनुमान प्रतिरूप है जिसका उपयोग
नकारात्मक दावे सिद्ध करने में होता है। मान लें आप दिखाना चाहते हैं कि कोई
दावा $p$ \emph{असत्य} है, अर्थात् $\lnot p$ दिखाना चाहते हैं। सबसे आशाजनक
रणनीति है (a) $p$ को सत्य मानना और (b) दिखाना कि यह मान्यता ऐसी बात तक ले
जाती है जिसे आप असत्य जानते हैं। ``असत्य ज्ञात बात'' ऐसा परिणाम हो सकता है
जो स्वयं $p$ से टकराता---उसका विरोध करता---है, अथवा विचाराधीन समग्र दावे की
किसी अन्य परिकल्पना से। उदाहरणतः ``यदि $q$, तो $\lnot p$'' के प्रमाण में
$q$ को सत्य मानना और उससे $\lnot p$ सिद्ध करना होता है। यदि $\lnot p$ को
विरोधाभास से सिद्ध करते हैं, तो इसका अर्थ $q$ के अतिरिक्त $p$ भी मानना है।
यदि $p$ से $\lnot q$ सिद्ध कर सकें, तो दिखा दिया कि मान्यता $p$ ऐसी बात तक
ले जाती है जो आपकी दूसरी मान्यता $q$ का विरोध करती है, क्योंकि $q$ और
$\lnot q$ दोनों सत्य नहीं हो सकते। निस्संदेह विरोधाभास के प्रमाण में
परिभाषाएँ खोलने के साथ अन्य अनुमान प्रतिरूप भी उपयोग करने होंगे। एक उदाहरण
देखें।

\begin{prop}
  यदि $A\subseteq B$ और $B=\emptyset$, तो $A$ के कोई !!{element} नहीं हैं।
\end{prop}

\begin{proof}
  मान लें $A\subseteq B$ और $B=\emptyset$। हम दिखाना चाहते हैं कि $A$ के
  कोई !!{element} नहीं हैं।
  \begin{quote}
    चूँकि यह सशर्त दावा है, हम पूर्ववर्ती मानते हैं और अनुवर्ती सिद्ध करना
    चाहते हैं। अनुवर्ती है: $A$ के कोई !!{element} नहीं हैं। इसे कुछ अधिक
    स्पष्ट कर सकते हैं: ऐसा नहीं है कि कोई $x\in A$ है।
  \end{quote}
  $A$ के कोई !!{element} नहीं हैं तभी जब ऐसा नहीं है कि कोई $x$ है जिसके
  लिए $x\in A$।
  \begin{quote}
    अतः हमने निर्धारित किया कि वास्तव में हमें नकारात्मक दावा $\lnot p$
    सिद्ध करना है, अर्थात्: ऐसा नहीं है कि कोई $x\in A$ है। विरोधाभास द्वारा
    प्रमाण के लिए संगत सकारात्मक दावा $p$, अर्थात् कोई $x\in A$ है, मानकर
    उससे विरोधाभास सिद्ध करना होगा। ``विरोधाभास के लिए मानें'' अथवा केवल
    ``अन्यथा मानें'' लिखकर और फिर मान्यता $p$ कहकर संकेत करते हैं कि
    विरोधाभास द्वारा प्रमाण कर रहे हैं।
  \end{quote}
  अन्यथा मानें: कोई $x\in A$ है।
  \begin{quote}
    अब यह नई मान्यता है जिसका उपयोग विरोधाभास प्राप्त करने में करेंगे।
    हमारी दो और मान्यताएँ हैं: $A\subseteq B$ और $B=\emptyset$। पहली से
    $x\in B$ मिलता है:
  \end{quote}
  चूँकि $A\subseteq B$, $x\in B$।
  \begin{quote}
    किंतु चूँकि $B=\emptyset$, $B$ का प्रत्येक !!{element} (जैसे $x$),
    $\emptyset$ का भी !!a{element} होना चाहिए।
  \end{quote}
  चूँकि $B=\emptyset$, $x\in\emptyset$। यह विरोधाभास है, क्योंकि परिभाषा
  से $\emptyset$ के कोई !!{element} नहीं हैं।
  \begin{quote}
    इससे प्रमाण पूरा हो चुका है: व्यवस्थित की गई मान्यताओं से हमें आवश्यक
    वस्तु (विरोधाभास) मिल गई और इसका अर्थ है कि सभी मान्यताएँ सत्य नहीं हो
    सकतीं। चूँकि पहली दो मान्यताएँ ($A\subseteq B$ और $B=\emptyset$)
    विवादित नहीं, अंतिम प्रस्तुत मान्यता (कोई $x\in A$ है) ही असत्य होनी
    चाहिए। किंतु पूर्ण विस्तार चाहें तो इसे स्पष्ट लिख सकते हैं।
  \end{quote}
  अतः हमारी मान्यता कि कोई $x\in A$ है, असत्य होनी चाहिए; फलतः विरोधाभास
  द्वारा प्रमाण से $A$ के कोई !!{element} नहीं हैं।
\end{proof}

प्रत्येक सकारात्मक दावा तुच्छतः किसी नकारात्मक दावे के तुल्य है: $p$ तभी
जब $\lnot\lnot p$। अतः विरोधाभास द्वारा प्रमाण का उपयोग सकारात्मक दावे
``अप्रत्यक्ष रूप से'' स्थापित करने में भी इस प्रकार हो सकता है: $p$ सिद्ध
करने के लिए उसे नकारात्मक दावा $\lnot\lnot p$ पढ़ें। यदि $\lnot p$ से
विरोधाभास सिद्ध कर सकें, तो विरोधाभास द्वारा प्रमाण से $\lnot\lnot p$ और
फलतः $p$ स्थापित किया है।

पिछले उदाहरण में हमारा लक्ष्य नकारात्मक दावा सिद्ध करना था, अर्थात् $A$ के
कोई !!{element} नहीं हैं, इसलिए विरोधाभास द्वारा प्रमाण के उद्देश्य से की
गई मान्यता (अर्थात् कोई $x\in A$ है) सकारात्मक दावा थी। उसने हमें काम करने
के लिए काल्पनिक $x\in A$ दिया, जिसके बारे में तब तक तर्क किया जब तक
$x\in\emptyset$ न मिला।

सकारात्मक दावा अप्रत्यक्ष रूप से सिद्ध करते समय विरोधाभास द्वारा प्रमाण के
लिए की जाने वाली मान्यता नकारात्मक होगी। किंतु बहुत बार सकारात्मक दावे को
सरलता से नकारात्मक और नकारात्मक दावे को सकारात्मक रूप में फिर सूत्रबद्ध कर
सकते हैं। यदि पिछले प्रमाण में नकारात्मक अनुवर्ती ``$A$ के कोई !!{element}
नहीं'' के बदले ``$A=\emptyset$'' सिद्ध करते, तो प्रमाण मूलतः समान होता।
($=$ की परिभाषा से ``$A=\emptyset$'' सामान्य दावा है, क्योंकि वह ``$A$ का
प्रत्येक !!{element}, $\emptyset$ का !!{element} है और इसके विपरीत'' में
खुलता है।) किंतु उसे नकारात्मक दावा ``ऐसा नहीं: कोई $x\in A$ है'' के तुल्य
सरलता से देखा जा सकता है।

अतः कभी-कभी $p$ को सीधे सिद्ध करने की अपेक्षा $\lnot p$ को मान्यता बनाकर
काम करना सरल है। जब सीधा प्रमाण उतना ही सरल या और सरल हो (जैसे अगले
उदाहरणों में), तब भी कुछ लोग अप्रत्यक्ष रूप से आगे बढ़ना पसंद करते हैं। यदि
दोहरा निषेध भ्रमित करे, तो किसी दावे के विरोधाभास द्वारा प्रमाण को
\emph{विपरीत} दावे से विरोधाभास के प्रमाण के रूप में सोचें। अतः $\lnot p$
का विरोधाभास द्वारा प्रमाण मान्यता $p$ से विरोधाभास का प्रमाण है; और $p$ का
विरोधाभास द्वारा प्रमाण $\lnot p$ से विरोधाभास का प्रमाण है।

\begin{prop}
$A \subseteq A \cup B$.
\end{prop}

\begin{proof}
  हम $A\subseteq A\cup B$ दिखाना चाहते हैं।
  \begin{quote}
    प्रथम दृष्टि में यह सकारात्मक दावा है: प्रत्येक $x\in A$, $A\cup B$ में
    भी है। इसका निषेध है: कोई $x\in A$, $\notin A\cup B$ है। अतः इस निषेधित
    दावे को मानकर और यह दिखाकर कि वह विरोधाभास तक ले जाता है, दावे को
    अप्रत्यक्ष रूप से सिद्ध कर सकते हैं।
    \end{quote}
  अन्यथा मानें, अर्थात् $A\nsubseteq A\cup B$।
  \begin{quote}
    हमारे पास $A\subseteq A\cup B$ की परिभाषा है: प्रत्येक $x\in A$,
    $\in A\cup B$ भी है। $A\nsubseteq A\cup B$ का अर्थ समझने के लिए खुली
    परिभाषा पर कुछ प्राथमिक तार्किक परिचालन करना होगा: यह असत्य है कि
    प्रत्येक $x\in A$, $\in A\cup B$ भी है, तभी जब कोई $x\in A$ ऐसा है जो
    $\notin C$। (यहाँ बहुत सावधान रहना चाहिए: अनेक विद्यार्थियों के
    विरोधाभास द्वारा प्रमाण के प्रयास इसलिए विफल होते हैं कि वे ``सभी $A$,
    $B$ हैं'' जैसे दावे के निषेध का गलत विश्लेषण करते हैं।) दूसरे शब्दों में
    $A\nsubseteq A\cup B$ तभी जब कोई $x$ है ऐसा कि $x\in A$ और
    $x\notin A\cup B$। इसके बाद सरल है।
  \end{quote}
  अतः कोई $x\in A$ है ऐसा कि $x\notin A\cup B$। $\cup$ की परिभाषा से
  $x\in A\cup B$ तभी जब $x\in A$ अथवा $x\in B$। चूँकि $x\in A$, हमारे
  पास $x\in A\cup B$ है। यह मान्यता $x\notin A\cup B$ का विरोध करता है।
\end{proof}

\begin{prob}
\emph{अप्रत्यक्ष रूप से} सिद्ध कीजिए कि $A\cap B\subseteq A$।
\end{prob}

\begin{prop}
यदि $A\subseteq B$ और $B\subseteq C$, तो $A\subseteq C$।
\end{prop}

\begin{proof}
  मान लें $A\subseteq B$ और $B\subseteq C$। हम $A\subseteq C$ दिखाना
  चाहते हैं।
  \begin{quote}
    अप्रत्यक्ष रूप से आगे बढ़ें: जिसे स्थापित करना चाहते हैं उसका निषेध
    मानते हैं।
  \end{quote}
  अन्यथा मानें, अर्थात् $A\nsubseteq C$।
  \begin{quote}
    पहले की तरह हम तर्क करते हैं कि $A\nsubseteq C$ तभी जब प्रत्येक
    $x\in A$, $\in C$ भी नहीं है, अर्थात् कोई $x\in A$, $\notin C$ है।
    चिंता न करें, अभ्यास के बाद ऐसे निषेध खोलने में बहुत सोचना नहीं पड़ेगा।
  \end{quote}
  दूसरे शब्दों में कोई $x$ है ऐसा कि $x\in A$ और $x\notin C$।
  \begin{quote}
    अब इसका उपयोग विरोधाभास तक पहुँचने में कर सकते हैं। निस्संदेह इसके लिए
    अन्य दो मान्यताएँ उपयोग करनी होंगी।
  \end{quote}
  चूँकि $A\subseteq B$, $x\in B$। चूँकि $B\subseteq C$, $x\in C$। किंतु
  यह $x\notin C$ का विरोध करता है।
\end{proof}

\begin{prop}
यदि $A\cup B=A\cap B$, तो $A=B$।
\end{prop}

\begin{proof}
  मान लें $A\cup B=A\cap B$। हम $A=B$ दिखाना चाहते हैं।
  \begin{quote}
    आरंभ अब नियमित है:
  \end{quote}
  विरोधाभास के लिए मान लें कि $A\neq B$।
  \begin{quote}
    विरोधाभास द्वारा प्रमाण के लिए हमारी मान्यता $A\neq B$ है। चूँकि
    $A=B$ तभी जब $A\subseteq B$ और $B\subseteq A$, हमें मिलता है कि
    $A\neq B$ तभी जब $A\nsubseteq B$ \emph{अथवा} $B\nsubseteq A$। (ध्यान
    दें कि निषेधों का परिचालन करते समय सावधानी कितनी महत्त्वपूर्ण है!{}) इस
    वियोजन से विरोधाभास सिद्ध करने के लिए स्थितियों द्वारा प्रमाण उपयोग करते
    और दिखाते हैं कि प्रत्येक स्थिति में विरोधाभास निष्पन्न होता है।
  \end{quote}
  $A\neq B$ तभी जब $A\nsubseteq B$ अथवा $B\nsubseteq A$। हम स्थितियों में
  भेद करते हैं।
  \begin{quote}
    पहली स्थिति में हम $A\nsubseteq B$ मानते हैं, अर्थात् किसी $x$ के लिए
    $x\in A$ किंतु $\notin B$। $A\cap B$ को उन !!{element}ों के रूप में
    परिभाषित किया गया है जो $A$ और $B$ में साझे हैं, अतः यदि कोई वस्तु उनमें
    से एक में नहीं, तो सर्वनिष्ठ में नहीं। $A\cup B$, $A$ के साथ $B$ है,
    इसलिए दोनों में से किसी में भी वस्तु सम्मिलन में है। इससे मिलता है कि
    $x\in A\cup B$ किंतु $x\notin A\cap B$, और फलतः
    $A\cap B\neq A\cup B$।
  \end{quote}
  
  स्थिति 1: $A\nsubseteq B$। तब किसी $x$ के लिए $x\in A$ किंतु
  $x\notin B$। चूँकि $x\notin B$, तब $x\notin A\cap B$। चूँकि $x\in A$,
  $x\in A\cup B$। अतः $A\cap B\neq A\cup B$, जो मान्यता
  $A\cap B=A\cup B$ का विरोध करता है।

  स्थिति 2: $B\nsubseteq A$। तब किसी $y$ के लिए $y\in B$ किंतु
  $y\notin A$। पहले की तरह हमारे पास $y\in A\cup B$ किंतु
  $y\notin A\cap B$, इसलिए $A\cap B\neq A\cup B$, जो फिर
  $A\cap B=A\cup B$ का विरोध करता है।
\end{proof}

\end{document}
